% =============================================================================
%  preamble.tex
%  Configuración general de la tesis. Cargado desde tesis.tex.
% =============================================================================

% --- Codificación y tipografía ------------------------------------------------
\usepackage[utf8]{inputenc}
\usepackage[T1]{fontenc}
\usepackage{lmodern}                % Latin Modern (mejor que Computer Modern por defecto)
\usepackage{textcomp}
\usepackage{microtype}              % Mejoras tipográficas finas

% --- Idioma -------------------------------------------------------------------
\usepackage[spanish,es-tabla,es-nodecimaldot,es-noindentfirst]{babel}
\usepackage{csquotes}               % Comillas correctas; requerido por biblatex
\selectlanguage{spanish}

% --- Geometría de página (estándar de tesis) ---------------------------------
\usepackage[letterpaper,
            inner=3.5cm,            % margen interior (encuadernación)
            outer=2.5cm,
            top=3cm,
            bottom=3cm,
            headheight=15pt,
            headsep=1cm,
            footskip=1.2cm]{geometry}

% --- Interlineado e indentación ---------------------------------------------
\usepackage{setspace}
\onehalfspacing                     % 1.5 de interlineado (estándar de tesis)
\setlength{\parindent}{1.25em}
\setlength{\parskip}{0.2em plus 0.05em minus 0.05em}

% --- Matemáticas -------------------------------------------------------------
\usepackage{amsmath,amssymb,amsthm,mathtools}
\usepackage{bm}

% --- Gráficos, color y figuras ----------------------------------------------
\usepackage{graphicx}
\graphicspath{{imagenes/}}
\usepackage{xcolor}
\definecolor{azulUNAM}{HTML}{003B71}
\definecolor{oroUNAM}{HTML}{C8A047}
\definecolor{grisCodigo}{HTML}{F4F4F4}
\usepackage{caption}
\captionsetup{font=small,labelfont=bf,labelsep=period,format=plain}
\usepackage{subcaption}
\usepackage{float}

% --- Tablas ------------------------------------------------------------------
\usepackage{booktabs}               % \toprule, \midrule, \bottomrule
\usepackage{tabularx}
\usepackage{multirow}
\usepackage{array}

% --- Listas ------------------------------------------------------------------
\usepackage{enumitem}
\setlist{itemsep=0.25em,topsep=0.3em}

% --- Encabezados y pies ------------------------------------------------------
\usepackage{fancyhdr}
\usepackage{emptypage}              % Páginas en blanco verdaderamente en blanco
\pagestyle{fancy}
\fancyhf{}
\fancyhead[LE,RO]{\thepage}
\fancyhead[RE]{\nouppercase{\leftmark}}     % Capítulo
\fancyhead[LO]{\nouppercase{\rightmark}}    % Sección
\renewcommand{\headrulewidth}{0.4pt}
\renewcommand{\footrulewidth}{0pt}

\fancypagestyle{plain}{%
  \fancyhf{}%
  \fancyfoot[C]{\thepage}%
  \renewcommand{\headrulewidth}{0pt}%
}

% --- Formato de capítulos y secciones ---------------------------------------
\usepackage{titlesec}
\titleformat{\chapter}[display]
  {\normalfont\huge\bfseries\color{azulUNAM}}
  {\filright\Large\MakeUppercase{\chaptertitlename}\ \thechapter}
  {1ex}
  {\titlerule\vspace{1ex}\filright}
\titlespacing*{\chapter}{0pt}{-20pt}{30pt}

\titleformat{\section}
  {\normalfont\Large\bfseries\color{azulUNAM!90!black}}{\thesection}{0.75em}{}
\titleformat{\subsection}
  {\normalfont\large\bfseries}{\thesubsection}{0.6em}{}
\titleformat{\subsubsection}
  {\normalfont\normalsize\bfseries\itshape}{\thesubsubsection}{0.5em}{}

% --- Código fuente (listings) ------------------------------------------------
\usepackage{listings}
\lstdefinestyle{python}{
  language=Python,
  basicstyle=\ttfamily\footnotesize,
  keywordstyle=\color{azulUNAM}\bfseries,
  commentstyle=\color{gray!80!black}\itshape,
  stringstyle=\color{oroUNAM!80!black},
  backgroundcolor=\color{grisCodigo},
  numbers=left,
  numberstyle=\tiny\color{gray},
  numbersep=8pt,
  frame=single,
  rulecolor=\color{gray!40},
  framerule=0.4pt,
  showstringspaces=false,
  breaklines=true,
  breakatwhitespace=true,
  tabsize=4,
  captionpos=b,
  literate={á}{{\'a}}1 {é}{{\'e}}1 {í}{{\'i}}1 {ó}{{\'o}}1 {ú}{{\'u}}1
           {Á}{{\'A}}1 {É}{{\'E}}1 {Í}{{\'I}}1 {Ó}{{\'O}}1 {Ú}{{\'U}}1
           {ñ}{{\~n}}1 {Ñ}{{\~N}}1 {¿}{{?`}}1 {¡}{{!`}}1
}
\lstset{style=python}
\renewcommand{\lstlistingname}{Código}
\renewcommand{\lstlistlistingname}{Índice de códigos}

% --- Entornos de teoremas (en español) --------------------------------------
\theoremstyle{definition}
\newtheorem{definicion}{Definición}[chapter]
\newtheorem{ejemplo}{Ejemplo}[chapter]

\theoremstyle{plain}
\newtheorem{teorema}{Teorema}[chapter]
\newtheorem{proposicion}[teorema]{Proposición}
\newtheorem{lema}[teorema]{Lema}
\newtheorem{corolario}[teorema]{Corolario}

\theoremstyle{remark}
\newtheorem{observacion}{Observación}[chapter]

% --- Bibliografía (biblatex) -------------------------------------------------
% Si tienes biber instalado, cambia "backend=bibtex" por "backend=biber".
\usepackage[backend=bibtex,
            style=ieee,
            sorting=nyt,
            block=space,
            bibencoding=utf8]{biblatex}
\addbibresource{referencias.bib}

% --- Hipervínculos (debe cargarse casi al final) ----------------------------
\usepackage[hidelinks,
            colorlinks=false,
            pdfstartview=Fit,
            pdfusetitle,
            bookmarksnumbered=true]{hyperref}

% --- Comandos personalizados -------------------------------------------------
\newcommand{\INS}{\textsc{ins}\xspace}
\newcommand{\PLN}{\textsc{pln}\xspace}
\newcommand{\NLP}{\textsc{nlp}\xspace}
\usepackage{xspace}

% Cita en bloque destacada
\newenvironment{citablock}
  {\begin{quote}\itshape\small}
  {\end{quote}}
